%% kapa: exam=1 round=30 year=2019 period=1 form=A subject=civil_law subject_ko=민법 seq=5 number=5
%% source: Q-Net 감정평가사 기출문제 (cst00309) artlSeq=5205205 file=extracted/1008/5205205/2019년 제30회 감정평가사 1차시험 문제지/2019년 제30회 감정평가사 1차시험 문제지/2019년도 제30회 감정평가사 1차 1교시 A형.hwp
%% conversion: kordoc 4.24.51 -> tools/gen1.py (auto); review=proofread-sonnet; flags=none
\begin{question}{5}
권리능력에 관한 설명으로 옳지 않은 것은? (다툼이 있으면 판례에 따름)
\begin{choices}
  \item 사람은 생존한 동안 권리와 의무의 주체가 된다.
  \item 사람이 권리능력을 상실하는 사유로는 사망이 유일하다.
  \item 수인(數人)이 동일한 위난으로 사망한 경우, 그들은 동시에 사망한 것으로 추정되므로 이 추정이 깨어지지 않는 한 그들 사이에는 상속이 일어나지 않는다.
  \item 의사의 과실로 태아가 사망한 경우, 태아의 부모는 태아의 의사에 대한 손해배상채권을 상속하여 행사할 수 있다.
  \item 인정사망에 의한 가족관계등록부에의 기재는 그 기재된 사망일에 사망한 것으로 추정하는 효력을 가진다.
\end{choices}
\end{question}
\answer{4}
